\begin{tabular}{lrllr}
\toprule
 & participação & tema & descricao & percepcao\_desfecho \\
\midrule
ESC\_INSE & 0.128 & Nível socioeconômico (INSE) & INSE da escola (1 = Nível I … 8 = Nível VIII) & False \\
DIR\_TX\_Q141\_\_FORA\_ESCALA & 0.063 & Merenda escolar & Quantas refeições são oferecidas nesta escola em relação ao tempo de permanência dos(as) estudantes? Para estudantes que permanecem mais que 4 e menos que 7 horas na escola: & False \\
DIR\_TX\_Q115 & 0.051 & Gestão democrática e colegiados & Existe Associação de Pais e Mestres – APM (ou caixa escolar) nesta escola? & False \\
DIR\_TX\_Q218 & 0.032 & Formação continuada percebida (equipe) & Educação especial. & False \\
DIR\_TX\_Q117 & 0.031 & Gestão democrática e colegiados & Há grêmio estudantil na sua escola? & False \\
DIR\_TX\_Q099 & 0.030 & Segurança — episódios de violência & Sobre os episódios listados abaixo, indique a frequência com que ocorreram neste ano, nesta escola: Discriminação (racial, gênero, orientação sexual, econômica/social, deficiência). & False \\
DIR\_TX\_Q044\_\_FORA\_ESCALA & 0.027 & Infraestrutura — Educação Infantil & Avalie, abaixo, as condições dos EQUIPAMENTOS da sua escola: Tanque de areia. & False \\
DIR\_TX\_Q170 & 0.024 & Composição de turmas e alocação de professores & Capacidade física da sala de aula. & False \\
DIR\_TX\_Q144 & 0.024 & Merenda escolar & Em relação à merenda escolar, como você avalia os seguintes aspectos: Os alimentos são de boa qualidade. & False \\
PROF\_TX\_Q023 & 0.024 & Formação continuada (curta duração) & Curso com carga horária total com mais de 180 e menos 360 horas & False \\
PROF\_TX\_Q137 & 0.024 & Clima escolar — segurança e violência & Roubo ou furto. & False \\
PROF\_TX\_Q021 & 0.023 & Formação continuada (curta duração) & Atividades formativas com menos de 20 horas & False \\
DIR\_TX\_Q004 & 0.023 & Caracterização da escola (etapas oferecidas) & Anos Finais do Ensino Fundamental. & False \\
DIR\_TX\_Q191 & 0.023 & Ações contra repetência e abandono & Neste ano, a escola realizou as seguintes ações para redução do ABANDONO ESCOLAR? Caso tenham sido realizadas, indique a eficácia das seguintes ações: Entrar em contato com os familiares/responsáveis do(a) estudante. & False \\
DIR\_TX\_Q014 & 0.023 & Formação e carreira do diretor & Qual é o MAIS ALTO nível de escolaridade que você concluiu? & False \\
PROF\_TX\_Q128 & 0.023 & Clima escolar — relação com estudantes & São assíduos(as). & False \\
PROF\_TX\_Q092 & 0.019 & Práticas pedagógicas em sala de aula & Buscar estratégias para melhorar a aprendizagem dos(as) estudantes com menor desempenho. & False \\
PROF\_TX\_Q009 & 0.018 & Hábitos culturais e lazer do professor & Acesso blogs, Youtube, redes sociais (Twitter, Instagram, Facebook etc.). & False \\
GEO\_ID\_AREA & 0.018 & Contexto geográfico & Capital ou interior & False \\
DIR\_TX\_Q121 & 0.017 & Parcerias externas & A escola desenvolve REGULARMENTE trabalhos em conjunto com: Serviços de saúde (postos de saúde). & False \\
DIR\_TX\_Q067 & 0.015 & Recursos financeiros e de pessoal & As famílias contribuíram com o trabalho pedagógico. & False \\
PROF\_TX\_Q077 & 0.015 & Crenças pedagógicas (repetência, avaliações externas) & A quantidade de avaliações externas (municipais, estaduais ou federais) é excessiva. & False \\
DIR\_TX\_Q043\_\_FORA\_ESCALA & 0.015 & Infraestrutura — Educação Infantil & Avalie, abaixo, as condições dos EQUIPAMENTOS da sua escola: Horta. & False \\
DIR\_TX\_Q079 & 0.014 & Interrupções do calendário & Falta de energia. & False \\
DIR\_TX\_Q001 & 0.014 & Caracterização da escola (etapas oferecidas) & Educação Infantil  -  Creche  (0 a 3 anos). & False \\
\bottomrule
\end{tabular}
